% Generated draft: requires mathematical and source review.
\documentclass[11pt]{article}
\usepackage[margin=25mm]{geometry}
\usepackage{fontspec}
\setmainfont{Latin Modern Roman}
\usepackage{amsmath,amssymb,mathtools,unicode-math}
\setmathfont{Latin Modern Math}
\usepackage{enumitem,needspace,xurl,hyperref}
\setlength{\parindent}{0pt}
\setlength{\parskip}{5pt}
\setlength{\emergencystretch}{3em}
\newcommand{\N}{\mathbb{N}}
\newcommand{\Z}{\mathbb{Z}}
\newcommand{\Q}{\mathbb{Q}}
\newcommand{\R}{\mathbb{R}}
\newcommand{\C}{\mathbb{C}}
\newcommand{\F}{\mathbb{F}}
\newcommand{\A}{\mathbb{A}}
\newcommand{\PP}{\mathbb{P}}
\newcommand{\E}{\mathbb{E}}
\newcommand{\eps}{\varepsilon}
\newcommand{\dd}{\,\mathrm{d}}
\begin{document}
\section*{UnsolvedMath ID: 30003620}
\textbf{Source problem code:} catalogue entry
\subsection{Definitions and mathematical statement}
For an integer \(d\geq 1\), define
\[
\mathcal P_d^{\ast}=\left\{\sum_{j=0}^{d}a_jX^j:\ a_j\in\{-1,0,1\}\text{ for every }j,\ (a_0,\ldots,a_d)\neq(0,\ldots,0)\right\}.
\]
Thus \(\mathcal P_d^{\ast}\) consists of the nonzero polynomials of degree at most \(d\) with the specified coefficients. Under the conventional interpretation that roots are complex numbers and that the zero polynomial is excluded, the proposed question is whether
\[
\exists s>0\ \forall d\in\Z_{\geq1}\ \forall P,Q\in\mathcal P_d^{\ast}\ \forall\alpha,\beta\in\C,
\qquad
\bigl(P(\alpha)=0\ \land\ Q(\beta)=0\ \land\ \alpha\neq\beta\bigr)
\Longrightarrow |\alpha-\beta|>s^d.
\]
The constant \(s\) must be independent of \(d,P,Q,\alpha,\beta\). The polynomials \(P\) and \(Q\) may coincide; the condition concerns distinct root values, without any simplicity or multiplicity assumption.

\textit{Unresolved source qualifications.} The supplied statement does not explicitly specify the root domain or exclude the zero polynomial. The displayed formulation makes the complex-root interpretation and nonzero-polynomial restriction explicit, subject to human confirmation. If the zero polynomial is admitted, its treatment must be specified before the source question has an unambiguous intended meaning. The supplied wording is interpreted as allowing roots of different polynomials; it does not explicitly say that both roots must belong to one polynomial. The positive-integer range for \(d\) is conventional; including \(d=0\) adds only a vacuous case under the nonzero-polynomial restriction.
\subsection{Short English statement}
Is there one positive constant \(s\) such that, for every positive integer \(d\), distinct complex roots drawn from any two nonzero polynomials of degree at most \(d\) with coefficients in \(\{-1,0,1\}\) are more than \(s^d\) apart? The complex-root domain, exclusion of the zero polynomial, and interpretation allowing different polynomials require confirmation from the source.
\subsection{Sources}
\begin{itemize}
\item Supplied local UnsolvedMath registry record, permanent ID \(30003620\), source problem code \texttt{catalogue entry}, titled ``Exponential Separation of Roots of Bounded Coefficient Polynomials,'' in the Oberwolfach Reports Open Problems collection. No individual proposer is supplied; the proposed year is recorded as 2017.
\item Supplied report citation: \emph{Arbeitsgemeinschaft: Additive Combinatorics, Entropy, and Fractal Geometry} (2018), Oberwolfach Reports, pp.~2847--2905, DOI \texttt{10.4171/owr/2017/47}. \url{https://doi.org/10.4171/owr/2017/47}. The supplied literature-source label is ``Oberwolfach Report 47 2017.'' The citation's 2018 date and the registry's 2017 year are retained without reconciliation.
\item Supplied literature reference labeled ``Breuillard Varju'': \url{https://arxiv.org/abs/1610.09154}. No paper title or further bibliographic details are supplied.
\item Literature assessment dated exactly \texttt{2026-08-21}: the supplied curation labels the problem open and states that the requested uniform exponential separation is not known and that current general lower bounds are weaker than exponential in the degree. This is supplied dated curation, not an independently verified assessment; the background expressly requests independent verification of time-sensitive status and cited claims.
\item The registry labels the statement \texttt{corrected\_verified} and reports that the official Oberwolfach report was checked directly, with definitions restored and source-dependent wording removed. This is a supplied verification claim, not verification performed for this draft.
\end{itemize}
\end{document}
